\subsection{An inverse trace need not be an invertible object}
The stack construction in \fid{F35} permits convergent reciprocal
traces. The next calculation identifies the additional obstruction
to interpreting such a reciprocal as a tensor inverse in the usual
derived Weil realization.

\begin{finding}[\fid{F38}: trace inversion and tensor inversion differ]\label{f:traceinverse}
Over $\F_q$, set
\[
 \mathcal Z=\A^1\times\mathbb G_m^2\times B\GL_2,
 \qquad \mathcal W=\mathbb P^1\times\mathcal Z.
\]
For every $Q=q^n$ their groupoid masses are
\[
 \#\mathcal Z(\F_Q)=\frac1{Q+1},\qquad
 \#\mathcal W(\F_Q)=1.
\]
The corresponding compact Euler traces converge absolutely.
Nevertheless $R\Gamma_c(\mathcal W)$ is not the unit: it has
nonzero cohomology in arbitrarily negative degrees, including
$H_c^{-1}(\mathcal W)=\Q_\ell(1)^2$.

In the usual unbounded derived category of Weil vector spaces,
with the algebraic tensor product,
\[
 H=R\Gamma(\mathbb P^1_{\bar{\F}_q},\Q_\ell)
   =\Q_\ell\oplus\Q_\ell(-1)[-2]
\]
has no tensor inverse. More generally, no tensor-invertible object
in this category has all-extension Euler trace
$D(q^n)=q^n+2+q^{-n}$.
Thus a strong monoidal realization with this specified image of
the projective-line object cannot extend through a tensor localization
making that object invertible. No faithfulness assumption is needed.
\end{finding}

\begin{proof}
Both stacks are smooth, of dimensions $-1$ and zero, and have
affine stabilizers. The torsor and column-count calculation in
\fid{F35} gives
\[
 \#\mathcal Z(\F_Q)
 =\frac{Q(Q-1)^2}{|\GL_2(\F_Q)|}
 =\frac1{Q+1}.
\]
Multiplication by $\#\mathbb P^1(\F_Q)=Q+1$ proves the mass
formula for $\mathcal W$.

Let $c_s=\lfloor s/2\rfloor+1$. Using the classifying-stack
calculation \cite[Example~7.4, p.~109]{SunStacks} and the compact
K\"unneth formula \cite[Theorem~11.0.14, p.~204]{LO2}, we obtain
\begin{equation}\label{eq:reciprocalstackgroups}
 H_c^j(\mathcal Z)=
 \bigoplus_{\substack{0\le b\le2,\ s\ge0\\j=-2-b-2s}}
 \Q_\ell(b+s+1)^{\oplus\binom2b c_s}.
\end{equation}
Indeed $\A^1$ contributes degree two and twist $-1$;
$\mathbb G_m^2$ contributes degree $4-b$ and twist $-(2-b)$;
and $B\GL_2$ contributes degree $-8-2s$ and twist $s+4$.
The sum of these degrees and twists is precisely the displayed
formula. The finite-field and finite-type hypotheses are the
same as those verified in \fid{F35}.

Put $K=R\Gamma_c(\mathcal Z)$. The projective bundle formula,
already used in \fid{F15}, gives the displayed decomposition of $H$.
Properness of $\mathbb P^1$ and K\"unneth give
\[
 R\Gamma_c(\mathcal W)=H\otimes K,\qquad
 H_c^j(\mathcal W)=H_c^j(\mathcal Z)
       \oplus H_c^{j-2}(\mathcal Z)(-1).
\]
Hence $H_c^0(\mathcal W)=\Q_\ell$ and
$H_c^{-1}(\mathcal W)=\Q_\ell(1)^2$.
Formula~\eqref{eq:reciprocalstackgroups} also makes the unbounded
negative support explicit. These are actual cohomology direct sums,
so their nonzero summands cannot cancel.

For $x=q^{-n}$ the compact Euler series of $K$ is
\[
 x\frac{(1-x)^2}{(1-x)(1-x^2)}
   =\frac{x}{1+x}
   =\sum_{j\ge0}(-1)^j x^{j+1}.
\]
Every positive real absolute eigenvalue moment converges: if $t>0$,
the ungrouped sum is
\[
 \sum_{b,s}\binom2b c_s q^{-t(b+s+1)}
 =\frac{q^{-t}(1+q^{-t})}{(1-q^{-t})^2}<\infty.
\]
Thus Sun's trace theorem applies, and this explicit convergence
also justifies all regrouping. Tensoring with the finite complex
$H$ preserves absolute convergence. Since its Euler trace is
$1+x^{-1}$, the Euler trace of $H\otimes K$ is exactly one.
After grouping, its only visible root is one with coefficient one;
the remaining actual cohomology cancels between degrees. The
uniqueness in \fid{F32} concerns that visible spectrum and is not
a faithfulness assertion for the trace.

For completeness, over a field every complex, including an
unbounded one, splits as its cohomology with zero differential
plus a contractible complex. Tensoring a contractible complex
remains contractible. With algebraic direct-sum totalization,
the resulting K\"unneth formula is
\[
 H^m(A\otimes B)=\bigoplus_{i+j=m}H^i(A)\otimes H^j(B).
\]
If $A\otimes B$ is the one-dimensional unit, both factors are
nonzero. Two independent total cohomology vectors in either
factor would give at least two in their tensor product.
Consequently each factor has one-dimensional total cohomology,
concentrated in a single degree. Forgetting the Weil action is
exact and conservative, so the same necessity holds for Weil
complexes. Truncation then identifies each factor with a
one-dimensional Weil representation in one degree. Conversely
such a representation, with its shift, has its dual with the
opposite shift as inverse.

In particular $H$ has no inverse. Directly, if $H^a(B)\ne0$,
then $H\otimes B$ has nonzero cohomology in both degrees $a$
and $a+2$. Strong monoidal functors preserve tensor inverses.
Therefore no strong monoidal realization sending a designated
object to this $H$ can factor through a localization in which
that object becomes invertible.

For the nonstandard normalization, under $v=Q^{\pm1/2}$ we have
\[
 D(Q)=(v+v^{-1})^2=Q^{-1}(Q+1)^2.
\]
The actual complexes
\[
 H_D=\Q_\ell(1)\otimes H^{\otimes2},\qquad
 K_D=\Q_\ell(-1)\otimes K^{\otimes2}
\]
have traces $D(Q)$ and $D(Q)^{-1}$, respectively. Their product
has trace one, but $H_D$ has total cohomology dimension four and
cannot have a tensor inverse. This conclusion does not depend
on choosing $H_D$: every invertible Weil complex has trace
$\pm\alpha^n$ for a single nonzero Frobenius scalar $\alpha$.
The minus sign cannot give the positive first and second values
of $D(q^n)$. The plus sign would require $D(q^2)=D(q)^2$, whereas
\[
 D(q)^2-D(q^2)=4(q+q^{-1}+1)>0.
\]
This proves the asserted all-extension obstruction.
\end{proof}

The statement permits formal Euler-class localizations, other
realizations, and completions with a different tensor product or
trace functional. It also permits multiplying an operator by the
numerical scalar $D(q)^{-1}$ at a fixed $q$. Scalar division does
not assert an object-level inverse or a reevaluation at every
$q^n$. The fixed realization and usual derived tensor product
are essential hypotheses of the obstruction.

% Separate paragraph for Part III.
\paragraph{\fid{F38} in Grothendieck's language.}
The projective-line decomposition and the stack's actual Tate
cohomology distinguish a positive object from its virtual Euler
class. A rational mass can invert an Euler trace while leaving
unbounded nonzero cohomology. Thus localization of a coefficient
ring or of a convergent trace algebra does not itself provide an
invertible geometric object. The split $\GL_2$ group scheme,
projective bundle formula, and modern stack six operations are
the precise frameworks used here. No archive leaf is claimed to
contain this construction or to identify it with a nonstandard
canonical basis. A different categorical localization requires a
specified new realization and its own comparison theorem.
